%!TEX root = ../dokumentation.tex
\chapter{Agility Discussion}

Im Projekt wird überwiegend ein klassisches Vorgehen nach dem Wasserfallmodell angewendet (Phasen Konzeption, Implementierung, Test, Rollout, Stabilisierung). Ein solcher Ansatz eignet sich, wenn Ziele und Umfang früh klar definierbar sind. Genau dies trifft hier zu. Der Business Case definiert einen klaren \ac{MOV} (Reduktion administrativer Aufwände, Nutzungsquote, Kompetenzkennzahlen) und bildet die Grundlage für die Projektplanung.

Das klassische Vorgehen ermöglicht Planbarkeit von Budget, Zeit und Ressourcen (zentral für die Charter). Stage-Gate-Entscheidungspunkte (z.\,B. nach Test oder Pilot) erlauben regelmäßige Prüfung der Zielbeiträge. Die klare Governance-Struktur (Verantwortlichkeiten, Meilensteine, Lieferobjekte) ist bei großer Organisation und vielen Stakeholdern von Bedeutung.

Agile Modelle wie \textit{Scrum} setzen auf iterative Entwicklung, kurze Feedbackzyklen und flexible Anpassung. Scrum arbeitet in \textit{Sprints} mit funktionsfähigen Teilergebnissen und hoher Anpassungsfähigkeit, was bei unklaren Anforderungen zu höherer Qualität und Nutzerakzeptanz führen kann sowie Nutzer durch Reviews einbindet.

Bei Standardsoftware wie der Lernplattform ist vollständiges Agil nur eingeschränkt geeignet, da viele Eigenschaften vorgegeben sind. Der Schwerpunkt liegt in Konfiguration, Integration, Datenmigration und organisatorischer Einführung, was sich besser phasenorientiert steuern lässt.

Agile Prinzipien können in Teilbereichen sinnvoll sein (iterative Konfigurations- und Testzyklen, Abstimmungen mit Fachbereichen, Verbesserungen in der Pilotphase). In Konfiguration und Pilotbetrieb kann feedbackorientiertes Vorgehen die Plattform besser an Nutzerbedürfnisse anpassen.

Für das Projekt ist ein hybrider Ansatz geeignet. Die übergeordnete Struktur folgt dem klassischen Phasenmodell (Planungssicherheit, Governance, Steuerbarkeit). Innerhalb von Phasen können agile Arbeitsweisen Flexibilität und Nutzerorientierung fördern. Das kombinierte Vorgehen verbindet Planbarkeit mit Anpassungsfähigkeit für eine erfolgreiche Einführung.